\begin{abstract}
Classical information theory is fundamentally limited by its local perspective,
analyzing pairwise interactions while ignoring the larger hierarchical
architecture of complex systems.
\emph{Structural Entropy} (SE) presents a paradigm shift: it extends Shannon
entropy to quantify information on a global scale and measures the uncertainty
embedded in a system's organizational hierarchy.
Since its introduction in 2016, SE has evolved from a principled measure for
community detection into a broadly applicable framework spanning graph learning,
reinforcement learning, bioinformatics, social-network forensics, and pattern
recognition---yet a systematic, unified synthesis of its theory, algorithms,
and applications has been lacking.

This survey provides a comprehensive decade-long overview of SE research.
We (i)~formalize its theoretical foundations, covering the seminal framework,
key extensions to directed, dynamic, and multi-relational graphs, and formal
connections to spectral entropy;
(ii)~survey the computational landscape, cataloguing heuristic, constrained,
and differentiable SE minimization algorithms together with their complexity
profiles;
(iii)~analyze SE-driven learning paradigms in graph pooling, structure
augmentation, and reinforcement learning;
(iv)~review cross-domain applications in bioinformatics, transportation,
geoscience, social networks, and pattern recognition;
and (v)~present a reproducibility benchmark comparing SE-based objectives
against competing graph-clustering methods.
We further discuss open research challenges and future directions, including
unified frameworks for heterogeneous and hypergraph structures, scalable
optimization, and the integration of SE into large language model pipelines.
A companion open-source library, \texttt{glass-jax}, accompanies this survey
as a reproducibility artifact.

\medskip\noindent\textbf{Keywords:} structural entropy, graph clustering,
hierarchical community detection, graph neural networks, reinforcement learning,
information theory.
\end{abstract}

